\subsection{Motivation and Background}

The task of generating reproducible research papers, particularly in machine learning, often faces the challenge of ensuring the reliability and consistency of experimental results. Recent advancements in multi-agent systems have shown promise in automating parts of the research process, but the reproducibility and robustness of these systems remain critical. This paper introduces PaperSwarm, an evidence-gated multi-agent system designed to address these challenges. The system leverages a novel approach where each agent is assigned a specific task, and their collective outputs are used to generate a coherent and reproducible paper.

\subsection{Problem Statement}

Consider a scenario where a machine learning model is trained on a dataset with \result{cl-features} features and \result{cl-train} training samples, and then evaluated on a separate test set of \result{cl-test} samples. In traditional settings, the performance of such a model can vary significantly due to factors like random initialization, hyperparameter tuning, and the inherent noise in the data. To mitigate these issues, we employ a multi-agent system that averages results over multiple random seeds, as indicated by \result{cl-seeds}.

\subsection{Proposed Solution}

PaperSwarm addresses the problem by integrating a multi-agent framework that collaborates to generate a paper. The system consists of several agents, each responsible for different tasks such as data preprocessing, model training, and result aggregation. For instance, one agent might focus on training a minimum-norm ordinary least squares (OLS) solution, while another handles ridge regression. The performance of these models is evaluated on the test set, and the results are averaged over \result{cl-seeds} to ensure robustness.

\subsection{Experimental Setup}

To evaluate the effectiveness of PaperSwarm, we conducted experiments on a synthetic dataset with \result{cl-features} features and \result{cl-train} training samples. The dataset was split into training and test sets, each containing \result{cl-test} samples. We compared the performance of two models: the minimum-norm OLS solution and ridge regression. The OLS model achieved a mean test mean squared error (MSE) of \result{cl-ols-mean} with a standard deviation of \result{cl-ols-std}, indicating high variance across different random seeds. In contrast, ridge regression, tuned to the penalty \result{cl-alpha}, achieved a mean test MSE of \result{cl-ridge-mean} with a much lower standard deviation of \result{cl-ridge-std}, suggesting more stable performance.

\subsection{Results}

The relative reduction in mean test MSE achieved by ridge regression over the OLS solution was \result{cl-improve} percent. This improvement underscores the benefits of incorporating regularization in the model training process. Furthermore, the irreducible noise floor, represented by \result{cl-noise}, serves as a benchmark for the true performance limits of the models. These results highlight the potential of PaperSwarm in generating reliable and reproducible research papers by leveraging the collective intelligence of a multi-agent system.

\subsection{Conclusion}

In summary, PaperSwarm demonstrates significant improvements in the reliability and robustness of experimental results. By employing a multi-agent system, the paper generation process becomes more transparent and reproducible, aligning with the principles of open science. Future work will explore the integration of additional agents and more complex models to further enhance the system's capabilities.
